%\section{Conclusão}

Neste trabalho realizou-se a identificação de um modelo matemático para o drone VTOL híbrido operando em modo multirrotor, voltado ao controle de decolagem, seguindo a metodologia de identificação no domínio do tempo M4V, com o método do erro de equação (EEM) fornecendo uma estimativa inicial, de baixo custo, e o método do erro de saída (OEM) refinando os parâmetros, {\color{blue}ajustando a saída simulada do modelo aos dados reais de voo}.

Para a obtenção destes dados, foram feitas várias campanhas de voo, dentre as quais algumas não foram bem sucedidas. {\color{blue}Houve voos em que o drone caiu e sofreu danos estruturais, reduzindo os dados aproveitáveis, voos em ambiente aberto nos quais o vento introduziu perturbações não modeladas, e variação de desempenho dos motores, que levou à substituição de um deles e exigiu recaracterizar o modelo de propulsão.} Essas dificuldades reforçaram a importância da seleção criteriosa dos segmentos dos dados, conforme os critérios de isolamento dos sinais de entrada, relação sinal-ruído e {\color{blue}largura do pulso do \textit{doublet}, executado próximo ao voo pairado}, e justificaram as escolhas metodológicas adotadas na identificação.

A modelagem partiu das equações de corpo rígido da dinâmica rotacional, acopladas a um modelo físico de propulsão, com dados obtidos em testes de bancada com os motores da aeronave. Um resultado central foi a comparação entre o modelo inicial, $\Theta_0$, montado a partir das inércias do modelo tridimensional e da caracterização em bancada dos motores, {\color{blue}e o modelo identificado por OEM}. Embora fisicamente consistente, o modelo inicial, por não ser ajustado a dados, não reproduz a dinâmica medida: os coeficientes de determinação $R^2$ das velocidades angulares $p$, $q$ e $r$ são fortemente negativos e a integração diverge. O modelo do OEM, ao contrário, reproduz bem as velocidades angulares ($R^2$ de {\color{blue}$0{,}71$, $0{,}73$ e $0{,}81$} em $p$, $q$ e $r$), com o resíduo concentrado na parcela de covariância, isto é, essencialmente aleatório, sem viés sistemático. {\color{blue}Conclui-se que o modelo construído apenas com medidas geométricas e ensaios de bancada é insuficiente para simulação e controle, e que a identificação com dados de voo é essencial para um modelo representativo.}

Sobre o modelo identificado projetou-se um controlador em arquitetura de cascata, com as malhas de atitude, velocidade e altitude ajustadas por alocação de polos sobre o modelo linearizado em torno do voo pairado. A malha fechada foi então simulada tanto no modelo linear quanto no modelo não linear completo de corpo rígido, em ambos os casos com a saturação real dos atuadores. A comparação entre as duas plantas delimitou a faixa de validade da linearização, {\color{blue}fiel para manobras dentro do envelope avaliado}, porém divergente em manobras agressivas e acopladas, nas quais os termos giroscópicos e a inclinação do vetor de empuxo, desprezados na linearização, passam a dominar a resposta. Esse mapeamento orienta o uso seguro do modelo linear no projeto de controle.

{\color{blue}A validação foi conduzida em janelas totalmente isoladas das de estimação, com manobras compostas nos três eixos, de modo que o desempenho reportado reflete generalização, e não sobreajuste.} Mais do que um conjunto de parâmetros, o trabalho entrega uma ferramenta de simulação realista para esta aeronave {\color{blue}em modo pairado}, capaz de servir de banco de provas para o projeto e a sintonia de controladores mais efetivos antes do voo real, reduzindo o risco e o custo dos ensaios em campo.

{\color{blue}Em síntese, o trabalho encadeou a modelagem dinâmica em modo multirrotor, a identificação dos parâmetros por EEM e OEM com dados reais de voo, o projeto do controlador em cascata e a comparação do controle nos modelos linear e não linear, constituindo uma base coerente do dado bruto de voo ao controlador simulado.}

{\color{blue}Diante das lacunas identificadas na Introdução, as contribuições originais deste trabalho podem ser enunciadas de forma objetiva. A primeira é a formulação com a matriz de inércia não diagonal, em que o produto $J_{xz}$ representa a assimetria herdada da fuselagem de asa fixa e acopla rolagem e guinada, tratamento pouco utilizado nos trabalhos da classe, que assumem multirrotores simétricos de matriz diagonal. A segunda é a estimação dos coeficientes efetivos de amortecimento angular $L_p$, $M_q$ e $N_r$ com incerteza inferior a $9\%$, coeficientes usualmente desconsiderados nos projetos de voo pairado de pequeno porte e reservados às aeronaves de maior porte. Os valores identificados resultam cerca de cinco, quatro e onze vezes maiores que os previstos pela teoria de influxo dos rotores, evidência quantitativa de que, no VTOL híbrido, além dos efeitos inerentes ao modo multirrotor existem acoplamentos da asa e da fuselagem imersas no fluxo induzido dos rotores, mecanismo distinto do amortecimento usual da asa, que depende da velocidade de translação. A terceira é a comparação em baixa velocidade do efeito dos termos aerodinâmicos da asa, que quantificou em até $6\%$ dos amortecimentos a contribuição dessas parcelas abaixo de $3$~m/s, sempre dentro do limite de Cramér-Rao, sustentando a forma reduzida do modelo. A quarta é o refinamento em voo dos coeficientes de bancada, com os empuxos corrigidos em cerca de $6\%$ e os contratorques entre $15\%$ e $29\%$, quantificando o desvio entre o ensaio estático e a operação real.}

{\color{blue}\section[{\color{blue}Limitações}]{{\color{blue}Limitações}}

O método e o modelo entregues têm um domínio de validade delimitado, e explicitá-lo é parte do resultado. A primeira fronteira é o envelope de velocidade: as parcelas aerodinâmicas da estrutura foram desprezadas no modelo final com respaldo da comparação da Seção~\ref{sec:aero-influencia}, válida para velocidades de translação abaixo de $3$~m/s. Acima desse envelope o modelo simplificado perde validade sem indicação interna, e a forma completa da Seção~\ref{sec:aero-estrutura} deve ser reativada. A segunda é espectral: o modelo foi verificado na banda das manobras, até cerca de $2$~Hz, e um controlador que exercite frequências superiores opera fora da região verificada, razão pela qual as bandas do projeto foram limitadas a $6$~rad/s na atitude e $2{,}5$~rad/s na guinada. A terceira é o canal de guinada: a baixa autoridade de contratorque limita a observabilidade, $J_z$ é determinado com o auxílio da restrição triangular ativa e $J_{xz}$ permanece indistinguível de zero pelos dados, contido em uma banda em torno do valor geométrico. Identificar esses parâmetros sem o auxílio de restrições exigiria manobras dedicadas de acoplamento entre rolagem e guinada, não contempladas na campanha. A quarta é estatística: os resíduos da estimação não são perfeitamente brancos nem gaussianos, comportamento esperado do OEM segundo a literatura, e por isso o limite de Cramér-Rao é lido como indicador relativo de confiança, e não como incerteza absoluta. Por fim, os parâmetros valem para a configuração de massa e centro de gravidade ensaiada, alterações de carga exigem nova identificação ou a robustez adicional discutida a seguir, e o controlador projetado foi avaliado somente em simulação.}

\section{Trabalhos Futuros}

O modelo de simulação não linear desenvolvido constitui a base natural para a identificação do modo de transição{\color{blue}, entre o voo pairado e o voo à frente. A via direta, de identificar a transição de uma só vez, não é bem posta, pois nesse regime não é possível dissociar as contribuições simultâneas dos rotores e da asa na sustentação e no arrasto. A alternativa parte dos dois extremos do envelope: o modo multirrotor, identificado neste trabalho no voo pairado, e o modo asa fixa, identificado com os rotores desligados. Com os dois extremos validados, os efeitos próprios da transição tornam-se isoláveis por diferença, e podem ser levantados em voos de decolagem seguidos de avanço horizontal com velocidade crescente até a proximidade do estol, confrontando a resposta medida com o que os dois modelos conhecidos explicam. Por isolar e validar a dinâmica do modo multirrotor, este trabalho entrega uma das metades dessa base, reduzindo a ambiguidade da estimação e separando as fontes de força e momento.}

\Needspace{4\baselineskip}
Como continuações diretas, destacam-se:
\begin{itemize}
    \item o embarque e o ensaio em voo do controlador projetado, etapa não realizada neste trabalho, cujo escopo se restringiu à avaliação em simulação sobre o modelo validado. Nessa etapa, a inclusão de ação integral com anti-\textit{windup} nas malhas de atitude e altitude, cujos ganhos se generalizam por alocação de polos de terceira ordem, confere robustez a variações de massa, à degradação da bateria e a perturbações de momento não representadas na simulação;
    \item o projeto de controladores mais efetivos, por exemplo com compensação direta dos termos de acoplamento ou com técnicas não lineares e robustas, tirando proveito da ferramenta de simulação aqui construída;
    \item a melhoria do canal de guinada, o mais fraco do modelo, por meio de manobras de excitação dedicadas{\color{blue}, incluindo as de acoplamento entre rolagem e guinada, capazes de tornar o produto de inércia observável}, da recaracterização do contratorque dos motores ou, principalmente, da substituição dos motores por outros com maior capacidade de contratorque{\color{blue}, o que também reduziria a perda de altitude associada à diferença de empuxo exigida pela manobra de guinada};
    \item o aprimoramento da instrumentação, com maior taxa de amostragem da IMU e a estimação dos erros sistemáticos do acelerômetro e do giroscópio como parâmetros do próprio OEM, para elevar a qualidade dos dados nas próximas campanhas;
    \item a estimação em tempo real, adicionando um observador ao controlador para atualizar o modelo durante a campanha de voo.
\end{itemize}
